\chapter{What nodwyn is}
\label{ch:what}

\section{The one-paragraph picture}

You type a message. A \textbf{front-end} (the app, a browser tab, the
terminal) hands it to a \textbf{session}. The session writes it to the
session's \textbf{log} and runs a \textbf{turn}: a Scheme procedure that asks
the \textbf{model} what to do. The model answers with text, or with
\textbf{tool calls} such as ``read \file{main.go}'' or ``run \code{go test}''.
Each tool call becomes a Scheme application, \code{(fs-read "main.go")}, and
is evaluated by the embedded Scheme interpreter. The interpreter has exactly
one way out to your computer: a \textbf{primitive}, a Go function in a table.
Every primitive application, whoever makes it, passes through one Go
function where the \textbf{guard} decides \emph{allow}, \emph{confirm} (ask
you) or \emph{deny}. Only then does the primitive run. Everything, including
the guard's decision, is appended to the log before anything else looks at
it, and the context the model sees next is rebuilt from that log.

That is the whole machine. The rest of the book is detail.

\begin{txt}
   you --> front-end --> session --> log
                            |
                            v
                  (nodwyn-turn request)            Scheme: yours to change
                            |
             +--------------+--------------+
             v                             v
       (model-turn) --> model        (run-tools calls)
                                           |
                                  (fs-read "main.go")
                                           |
                                   +-------v-------+
                                   |     GUARD     |   allow / confirm / deny
                                   +-------+-------+
                                           v
                              a Go primitive: the file, the program
\end{txt}

\section{Three ideas that explain everything}

\subsection{One door}

A program that wants to protect you usually sprinkles checks where it
thinks danger lives: here before deleting, there before sending. \nodwyn{}
does the opposite. There is \emph{one} door between the Scheme world and your
computer, the primitive dispatch, and the guard stands in it. A call from
the model, from your own \code{init.scm}, from the loop's strategy, from you
typing at the prompt, even a call made from inside \code{map} or
\code{sort}: all of them cross the same door. The test suite records every
crossing to prove there is no other way out.

Because the door is single, the guard can be simple and strict. Its
checkers look at \emph{what an argument is} (a path, a destination, a
payload, a credential) rather than guessing from the text, and the strictest
verdict wins: one \emph{deny} beats any number of \emph{allow}s.

\subsection{The log is the truth}

The model's context window is not memory; it is a \emph{view}. Every
request is rebuilt from the session log, an append-only file with one JSON
line per event, written and flushed before anything else happens. That buys
three things you will use as an expert:

\begin{itemize}
  \item \textbf{Replay.} \code{nodwyn replay SESSION} prints the whole
    conversation from the file alone, and \code{nodwyn replay SESSION -{}-window
    SEQ} rebuilds, byte for byte, the exact request the model saw at that
    moment, and checks its hash.
  \item \textbf{Audit.} Every guard decision that matters is a line: every
    confirm and deny, and every allow of a tool or of anything that touches
    your computer or the network (pure builtins like \code{+} are allowed
    silently). Each says who asked (you, the model, or the loop), what,
    which checker decided, why, and what you answered.
  \item \textbf{Freedom to fold.} Because nothing is ever lost from the log,
    a strategy can hide old tool output from the model, summarise a long
    conversation or run a sub-agent in a private context, and still keep
    everything for you.
\end{itemize}

\subsection{Credentials, not content}

Most agents try to protect secrets by recognising them in text. \nodwyn{}
turns the problem around. A secret you give it (an API key, say) is a
\textbf{credential}: you type it once into a masked prompt, it goes into an
encrypted keychain, and Scheme only ever holds an opaque \textbf{handle},
\code{\#<cred:openai-key:3f9a2c1e7b4d5a60>}, that no builtin can turn
into text. A
\textbf{binding} says where that credential may go: which host, which HTTP
methods, which position in the request (a header, a query parameter). The
key is injected at the last moment, on the way out, and only if the binding
table allows that exact destination.

Recognising secret-looking text still happens, as a second line:
\textbf{redaction} replaces recognised credentials with
\code{[REDACTED class]} in everything that enters the log or the model's
context, and a write whose content looks like a key needs your explicit
yes. But the first line is structural: the model never sees a key, so it
cannot leak one.

\section{Who is asking?}

Every primitive call carries an \textbf{origin}, and the guard treats them
differently:

\begin{center}
\begin{tabularx}{0.95\linewidth}{lX}
\toprule
\textbf{origin} & \textbf{who} \\
\midrule
\code{user} & you: \code{:eval} at the prompt, your \code{init.scm}, the REPL \\
\code{model} & the model's tool calls, and anything a tool it called does in turn \\
\code{harness} & the loop's own strategy: the planner, compaction, sub-agent plumbing \\
\bottomrule
\end{tabularx}
\end{center}

Some primitives are \emph{user-only}. Registering a guard check, writing a
policy file, registering a credential, removing provenance from a value:
the model can never do these, and neither can the loop's strategy. You will
meet each of them later.

\section{The promise, precisely}

It is worth being exact about what the guard promises, because an expert
knows where the edges are.

\begin{itemize}
  \item \textbf{Everything outside the folder is refused.} Files are pinned
    to the task's folder (plus a private scratch directory and any folders
    you \emph{mount} into a workspace). Symbolic links that point outside do
    not escape.
  \item \textbf{Credential-looking files are refused}: \file{.env} and \file{.env.*},
    \file{*.pem}, \file{id\_*}, \file{\textasciitilde/.ssh},
    \file{\textasciitilde/.aws}, \nodwyn's own config and data directories
    wherever \code{XDG\_CONFIG\_HOME} and \code{XDG\_DATA\_HOME} put them
    (all but the session's scratch directory, which lives in the data
    directory), and more.
  \item \textbf{A write in the folder can be undone.} By default it goes
    ahead without asking: its diff is logged and shown, and the file's
    previous content is kept, so it can be put back. A few paths always
    ask, every time: \file{.nodwyn/}, \file{.git/}, \file{.github/},
    \file{.gitea/}, \file{scripts/} and the \file{Makefile}, because what is
    written there later runs with your authority. Content that looks like a
    key asks too.
  \item \textbf{A command runs only when something says yes}: a standing
    approval, the judge, which is asked first by default, or you, when the
    judge refuses or there is none.
  \item \textbf{No key leaves unless a binding says where it may go}, and
    content that came from your files or from a command's output is refused
    on its way to a destination no binding names.
\end{itemize}

And what it does \emph{not} promise: that a command you or the judge
approve is safe. If you say yes to \code{bash -c "rm -rf build"}, it runs;
if you say yes to \code{curl}, it reaches the network with your access (the
guard flags network tools by name so you notice). The guard decides who
answers; it does not make the answers wise. Read the reason lines it shows
you. They are written to tell you what a yes would do. If you want every
command to be yours to answer, Chapter~\ref{ch:primitives} shows the
setting.

\begin{tip}
The fastest way to become good with \nodwyn{} is to read its log. Every
surprise (``why did it ask me that?'', ``why didn't it ask?'') has a
\code{guard.decision} line with a checker name and a reason. Chapter~\ref{ch:log}
shows how.
\end{tip}

\section{Why Scheme?}

Because the interesting parts of an agent are \emph{policy}: when to ask the
model again, what to do with a failed test, when to plan, what to forget.
Policy changes from project to project and week to week. \nodwyn{} keeps the
dangerous parts (files, processes, network, the keychain, the log) in Go,
small and fixed, and puts the policy in Scheme, where you can change it in
a file in your repository without rebuilding anything.

Scheme is a good fit for that job. It is tiny (you will learn all of the
dialect in Part~II), it treats code and data the same way, and a loop
written as a recursive function runs in constant space. It also cannot hurt
you by accident: the language has no file or network access of its own, so
the \emph{only} effects a Scheme program can have are primitive calls, and
every one of those meets the guard.

One more guarantee: loading Scheme can grow the table of primitives and
define new procedures, but it cannot change how Scheme itself evaluates.
\code{eval} and \code{apply} are two fixed Go functions, and the kernel has
a fingerprint the test suite checks before and after loading arbitrary
code. Your \code{init.scm} can rewire the loop. It cannot rewire the door.
